\documentclass[11pt]{article}

\usepackage[a4paper,margin=0.85in]{geometry}
\usepackage{graphicx}
\usepackage{booktabs}
\usepackage{longtable}
\usepackage{array}
\usepackage{hyperref}
\usepackage{xcolor}
\usepackage{enumitem}
\usepackage{fancyvrb}
\usepackage{float}

\hypersetup{
  colorlinks=true,
  linkcolor=blue!60!black,
  urlcolor=blue!60!black
}

\newcommand{\code}[1]{\texttt{\detokenize{#1}}}
\newcommand{\rootplot}[2]{
  \begin{figure}[H]
    \centering
    \includegraphics[width=0.96\linewidth]{../../../../#1}
    \caption{#2}
  \end{figure}
}
\newcommand{\edaplot}[2]{
  \begin{figure}[H]
    \centering
    \includegraphics[width=0.96\linewidth]{ve_vev_eda_output/charts/#1}
    \caption{#2}
  \end{figure}
}
\newcommand{\smallplot}[1]{
  \includegraphics[width=0.48\linewidth]{ve_vev_eda_output/charts/#1}
}

\title{Round 3 Velvetfruit and VEV Voucher EDA Synthesis}
\author{Prosperity 4 Trading Research}
\date{April 24, 2026}

\begin{document}
\maketitle

\section*{Scope And Source Files}

This document synthesizes two exploratory analyses:

\begin{itemize}
  \item \code{phase2/round3/algo/EDA.ipynb}, the first-pass dataset and order-book structure analysis.
  \item \code{phase2/round3/algo/ve_vev_eda.py}, the longer Velvetfruit and VEV voucher EDA that generated \code{ve_vev_eda_output/charts} and \code{ve_vev_eda_output/tables}.
\end{itemize}

The products in scope are \code{VELVETFRUIT_EXTRACT} and the ten call-like voucher products \code{VEV_4000}, \code{VEV_4500}, \code{VEV_5000}, \code{VEV_5100}, \code{VEV_5200}, \code{VEV_5300}, \code{VEV_5400}, \code{VEV_5500}, \code{VEV_6000}, and \code{VEV_6500}. The full Round 3 data also contains \code{HYDROGEL_PACK}; both EDAs either filter it out or explicitly skip it for the option-complex analysis.

\section*{Executive Findings}

\begin{itemize}
  \item The full combined price dataset has 360,000 rows: 12 products, 30,000 book snapshots each, split into 3 historical days and 10,000 timestamps per product per day.
  \item The notebook filters out \code{HYDROGEL_PACK} before computing product frequencies, so its printed denominator is 330,000 rather than the full 360,000.
  \item The target product trade counts match across the notebook and the longer EDA: \code{VELVETFRUIT_EXTRACT} has 1,372 trades; the most active vouchers are \code{VEV_4000}, \code{VEV_6000}, \code{VEV_6500}, \code{VEV_5500}, and \code{VEV_5400}.
  \item \code{VELVETFRUIT_EXTRACT} has a very stable top-of-book spread around 5 ticks, with bid and ask offsets centered around -2.5 and +2.5 from mid.
  \item Underlying one-tick log returns are non-normal and negatively autocorrelated at lag 1, but ten-tick returns look much closer to normal.
  \item The underlying has rolling annualized realized volatility around 0.41, while the ATM IV proxy sits around 0.235, producing a persistent IV-minus-RV gap near -0.178.
  \item The strongest gamma-scalp proxy appears in the near-ATM complex: \code{VEV_5200}, \code{VEV_5300}, \code{VEV_5400}, and \code{VEV_5100}.
  \item Deep OTM vouchers \code{VEV_6000} and \code{VEV_6500} are pinned at bid 0 / ask 1 / mid 0.5; they create large percentage spreads and fragile implied-volatility signals.
  \item Lead-lag correlations peak at lag 0 for all active vouchers. This argues against simple delayed cross-product arbitrage and toward delta/gamma fair-value modeling.
  \item No-arbitrage flags are concentrated in lower-bound violations for deep ITM \code{VEV_4000} and \code{VEV_4500}; these are likely driven by integer prices, near-zero extrinsic value, and midpoint conventions, not automatically executable arbitrage.
\end{itemize}

\section{Dataset Structure From \code{EDA.ipynb}}

\subsection*{Raw Shape And Product Coverage}

\begin{itemize}
  \item The full combined prices file is \code{phase2/round3/data/prices_round_3_combined.csv}.
  \item The full combined trades file is \code{phase2/round3/data/trades_round_3_combined.csv}.
  \item The price schema is:
  \code{day}, \code{timestamp}, \code{product}, three bid prices, three bid volumes, three ask prices, three ask volumes, \code{mid_price}, and \code{profit_and_loss}.
  \item The trade schema contains \code{timestamp}, \code{buyer}, \code{seller}, \code{symbol}, \code{currency}, \code{price}, \code{quantity}, and \code{day}.
  \item Every product appears at every 100-timestamp grid point from 0 to 999900 on each historical day.
\end{itemize}

\begin{Verbatim}[fontsize=\scriptsize]
Full combined prices:
HYDROGEL_PACK           30000
VELVETFRUIT_EXTRACT     30000
VEV_4000                30000
VEV_4500                30000
VEV_5000                30000
VEV_5100                30000
VEV_5200                30000
VEV_5300                30000
VEV_5400                30000
VEV_5500                30000
VEV_6000                30000
VEV_6500                30000

Full combined trades:
HYDROGEL_PACK           1010
VELVETFRUIT_EXTRACT     1372
VEV_4000                 464
VEV_4500                   1
VEV_5000                   1
VEV_5100                   1
VEV_5200                  18
VEV_5300                 121
VEV_5400                 225
VEV_5500                 267
VEV_6000                 284
VEV_6500                 284
\end{Verbatim}

\subsection*{Notebook Product-Frequency Print}

\begin{itemize}
  \item The notebook prints 30,000 rows per target product over 330,000 filtered rows.
  \item This is internally consistent after filtering, but it is not the full dataset denominator because \code{HYDROGEL_PACK} is absent from the filtered target list.
\end{itemize}

\begin{Verbatim}[fontsize=\scriptsize]
1. PRODUCT APPEARANCE FREQUENCIES IN ORDER BOOK
VEV_4000                 : 30000/330000 (09.09%)
VEV_4500                 : 30000/330000 (09.09%)
VEV_5000                 : 30000/330000 (09.09%)
VEV_5100                 : 30000/330000 (09.09%)
VEV_5200                 : 30000/330000 (09.09%)
VEV_5300                 : 30000/330000 (09.09%)
VEV_5400                 : 30000/330000 (09.09%)
VEV_5500                 : 30000/330000 (09.09%)
VEV_6000                 : 30000/330000 (09.09%)
VEV_6500                 : 30000/330000 (09.09%)
VELVETFRUIT_EXTRACT      : 30000/330000 (09.09%)
\end{Verbatim}

\rootplot{lob_prices_grid.png}{Notebook plot: order-book price distributions by product and layer.}
\rootplot{lob_volumes_grid.png}{Notebook plot: order-book volume distributions by product and layer.}

\subsection*{Order Book Structure Discoveries}

\begin{itemize}
  \item \code{VELVETFRUIT_EXTRACT} mid is centered near 5250 with top-of-book bid and ask means near 5247.6 and 5252.6.
  \item \code{VEV_4000} and \code{VEV_4500} behave like deep ITM calls: high absolute prices, large spreads, tiny extrinsic value.
  \item \code{VEV_5200} and \code{VEV_5300} are the most natural ATM candidates because spot is near 5250.
  \item \code{VEV_6000} and \code{VEV_6500} are pinned at bid 0, ask 1, mid 0.5 for all timestamps.
  \item Top-level voucher volumes are stable and integer quoted; the more OTM vouchers have increasingly shallow deeper-book availability.
  \item Average tick distance between order-book appearances is exactly 1 for all target products, meaning every filtered product has a quote at every timestamp grid point.
\end{itemize}

\begin{Verbatim}[fontsize=\scriptsize]
Mean L1 view from notebook:
Product                  bid_price_1   bid_volume_1   ask_price_1   ask_volume_1
VEV_4000                     1239.70          10.92       1260.52          10.92
VEV_4500                      742.18           8.95        758.04           8.95
VEV_5000                      252.00          15.33        258.04          15.48
VEV_5100                      164.66          19.25        168.95          19.35
VEV_5200                       94.10          22.56         96.99          22.51
VEV_5300                       45.71          20.28         47.81          20.24
VEV_5400                       15.26          21.76         16.64          21.72
VEV_5500                        6.07          22.18          7.22          22.17
VEV_6000                        0.00          22.48          1.00          22.48
VEV_6500                        0.00          15.48          1.00          15.48
VELVETFRUIT_EXTRACT          5247.60          37.83       5252.59          37.80
\end{Verbatim}

\rootplot{advanced_metrics_grid.png}{Notebook plot: advanced order-book metric distributions.}
\rootplot{tick_distances.png}{Notebook plot: tick distance between order-book appearances.}

\subsection*{Notebook Trade And Depletion Findings}

\begin{itemize}
  \item \code{VELVETFRUIT_EXTRACT} trades far more frequently than any voucher.
  \item Among vouchers, \code{VEV_4000} trades often but with tiny size around 2; near-ATM and OTM active vouchers trade in sizes around 3.5.
  \item \code{VEV_4500}, \code{VEV_5000}, and \code{VEV_5100} only have one trade each, all on day 1.
  \item \code{VEV_5200} is sparse, with 18 total trades across three days.
  \item The notebook's depletion scan finds 188 underlying events consuming at least layer 1, and only one event consuming at least layer 2/3: \code{VEV_6500} on day 0 at timestamp 640300.
  \item No empty L1, one-sided L1, posted-L2-without-L1, or posted-L3-without-L2 anomalies are found after the notebook fills missing volumes with zero.
\end{itemize}

\begin{Verbatim}[fontsize=\scriptsize]
Trades Count per Product per Day:
day       VEV_4000 VEV_4500 VEV_5000 VEV_5100 VEV_5200 VEV_5300 VEV_5400 VEV_5500 VEV_6000 VEV_6500 VELVETFRUIT_EXTRACT
0              172        0        0        0        3       37       64       81       91       91                 445
1              164        1        1        1        7       39       81       92       98       98                 450
2              128        0        0        0        8       45       80       94       95       95                 477

LOB depletion:
Instances consuming AT LEAST Layer 1:
VEV_5200                   7
VEV_5300                  19
VEV_5400                   5
VEV_5500                   5
VEV_6500                   1
VELVETFRUIT_EXTRACT      188

Instances consuming AT LEAST Layer 2 and Layer 3:
day=0, timestamp=640300, product=VEV_6500, total_trade_qty=5, avg_trade_price=0.0
\end{Verbatim}

\rootplot{trades_distributions.png}{Notebook plot: trade price and trade quantity distributions.}

\section{Data Integrity From The Longer EDA}

\begin{itemize}
  \item The longer EDA uses per-day files in \code{phase2/round3/algo/data}; these match the combined target-product counts from the notebook.
  \item Every target product has 10,000 book snapshots per day and no timestamp gaps.
  \item No crossed books, missing mids, duplicate timestamps, or trades without same-timestamp book rows are reported.
  \item Six \code{VELVETFRUIT_EXTRACT} trades are outside the visible top bid/ask across all days; this should be interpreted as hidden depth or timing artifact before being treated as edge.
  \item Deeper book availability varies sharply: \code{VEV_4000} and \code{VEV_4500} almost always have deeper levels, while \code{VEV_6000} and \code{VEV_6500} are L1-only.
\end{itemize}

\begin{Verbatim}[fontsize=\scriptsize]
Data integrity highlights:
VELVETFRUIT_EXTRACT: 10000 book snapshots per day, trades 445/450/477, outside top book 1/3/2.
VEV_4000:            10000 book snapshots per day, trades 172/164/128, no top-book anomalies.
VEV_5200:            10000 book snapshots per day, trades 3/7/8, no top-book anomalies.
VEV_6000:            10000 book snapshots per day, trades 91/98/95, L1-only book.
VEV_6500:            10000 book snapshots per day, trades 91/98/95, L1-only book.
\end{Verbatim}

\edaplot{data_integrity_trade_counts.png}{Long EDA: trade counts by product and day.}
\edaplot{data_integrity_depth_heatmap.png}{Long EDA: deeper-book availability by product and day.}

\section{Underlying: \code{VELVETFRUIT_EXTRACT}}

\subsection*{Price, Return, And Volatility Findings}

\begin{itemize}
  \item Mid price oscillates around 5250 with visible intraday trends and day-to-day continuation.
  \item Tick-level price-change standard deviation is about 1.12 to 1.14 price units.
  \item Day-level drift is small but positive on days 1 and 2: about +20.5 and +28.0 ticks per day, versus about -6.0 on day 0.
  \item Around 24.5\% of tick-level mid changes are exactly zero; positive and negative changes are otherwise balanced.
  \item One-tick log returns reject normality very strongly; ten-tick returns have much weaker non-normality and much higher Jarque-Bera p-values.
  \item Rolling realized volatility is stable: RV500 is near 0.410 to 0.414 by day.
  \item ACF of mid price is near 0.997 at lag 1, reflecting price-level persistence; ACF of one-tick returns is -0.159 at lag 1, reflecting bid/ask bounce or short-horizon mean reversion.
\end{itemize}

\begin{Verbatim}[fontsize=\scriptsize]
Return stats:
day  horizon  mean          std        skew      kurtosis   JB p-value
0    1       -1.14e-07     0.000214   -0.0397    0.4517    9.26e-20
1    1        3.90e-07     0.000216   -0.0307    0.3246    1.33e-10
2    1        5.30e-07     0.000217   -0.0294    0.2828    2.82e-08
pooled 10     3.08e-06     0.000576   -0.0024   -0.0073    0.9535

Stationarity diagnostics:
ADF p-value = 4.71e-05
KPSS p-value = 0.01
Variance ratios: VR(2)=0.841, VR(5)=0.748, VR(10)=0.715, VR(20)=0.705
Hurst R/S = 0.999
Script classification = trending
\end{Verbatim}

\edaplot{velvetfruit_mid_timeseries.png}{Underlying mid price across historical days.}
\edaplot{velvetfruit_realized_vol_day0.png}{Underlying rolling realized volatility on day 0.}
\edaplot{velvetfruit_realized_vol_day1.png}{Underlying rolling realized volatility on day 1.}
\edaplot{velvetfruit_realized_vol_day2.png}{Underlying rolling realized volatility on day 2.}
\edaplot{velvetfruit_return_hist_pooled_h1.png}{Pooled one-tick log returns.}
\edaplot{velvetfruit_return_hist_pooled_h10.png}{Pooled ten-tick log returns.}
\edaplot{velvetfruit_acf.png}{Underlying price and return autocorrelation.}

\subsection*{Liquidity, Flow, And Execution Findings}

\begin{itemize}
  \item The underlying spread is highly stable: mean 4.988, median 5, p95 6.
  \item Bid and ask quote offsets are symmetric around mid: -2.494 and +2.494 on average.
  \item Top L1 depth averages about 37.8 units per side, and L2/L3 average around 40 per side when present.
  \item Trade arrival is close to Poisson-like by 1000-timestamp buckets; dispersion indices are between 0.86 and 1.00.
  \item Trades are classified as 781 buys and 591 sells using trade price versus mid.
  \item Buy trades are followed by positive forward log returns at horizons 100, 500, and 1000; sell trades are weaker and only clearly negative at 100 and 1000.
  \item Maker fills are buy-flow biased: trades at ask are 55.3\% to 58.7\% of underlying prints by day.
  \item Crossing the spread costs about 2.494 ticks versus mid per taker trade, so naive taker hedging needs substantial option edge or urgent risk control.
\end{itemize}

\begin{Verbatim}[fontsize=\scriptsize]
Spread summary:
day  mean    median  mode  p95
0    4.9935  5.0     5     6.0
1    4.9846  5.0     5     6.0
2    4.9863  5.0     5     6.0

Flow toxicity:
horizon  side  mean_forward_return  count
100      buy       0.000079          781
100      sell     -0.000023          591
500      buy       0.000089          781
500      sell      0.000002          591
1000     buy       0.000075          781
1000     sell     -0.000030          591

Maker fills:
day  ask-side fills  bid-side fills  ask-fill fraction
0    252             193             0.566
1    249             201             0.553
2    280             197             0.587
\end{Verbatim}

\edaplot{velvetfruit_spread_hist.png}{Underlying top-of-book spread distribution.}
\edaplot{velvetfruit_spread_vs_rv.png}{Rolling spread versus rolling realized volatility.}
\edaplot{velvetfruit_trade_minus_mid_by_side.png}{Underlying trade price minus mid by inferred side.}
\edaplot{velvetfruit_trade_interarrival_hist.png}{Underlying trade interarrival distribution.}
\edaplot{velvetfruit_intraday_trade_direction.png}{Underlying intraday trade count and buy fraction.}
\edaplot{velvetfruit_quote_rounding.png}{Underlying quote rounding.}
\edaplot{velvetfruit_depth_profile.png}{Underlying depth profile.}

\section{Voucher Microstructure And Moneyness}

\begin{itemize}
  \item \code{VEV_4000}, \code{VEV_4500}, and \code{VEV_5000} are ITM across all days, although \code{VEV_5000} is closer to the money.
  \item \code{VEV_5100} transitions from mostly near-ATM to more ITM by day 2.
  \item \code{VEV_5200} and \code{VEV_5300} are the cleanest near-ATM strikes.
  \item \code{VEV_5400} is near-ATM but partly OTM, especially early.
  \item \code{VEV_5500}, \code{VEV_6000}, and \code{VEV_6500} are OTM across all days.
  \item Absolute spreads decline as strike rises, but spread as a percentage of mid explodes in OTM floor-price vouchers.
  \item \code{VEV_6000} and \code{VEV_6500} have spread-to-mid equal to 2.0 because mid is 0.5 and spread is 1.
  \item \code{VEV_4000} has many trades but wide absolute spreads around 21 ticks.
  \item \code{VEV_5300}, \code{VEV_5400}, and \code{VEV_5500} are more relevant for active option market-making because they combine nonzero trade activity with lower absolute spreads.
\end{itemize}

\begin{Verbatim}[fontsize=\scriptsize]
Voucher liquidity by day:
VEV_4000: mean spread about 20.8, mean mid 1246-1255, trades/day 172/164/128.
VEV_5200: mean spread about 2.86-2.93, mean mid 94-97, trades/day 3/7/8.
VEV_5300: mean spread about 2.05-2.16, mean mid 44-49, trades/day 37/39/45.
VEV_5400: mean spread about 1.33-1.43, mean mid 13.7-18.5, trades/day 64/81/80.
VEV_5500: mean spread about 1.12-1.18, mean mid 5.3-8.1, trades/day 81/92/94.
VEV_6000 and VEV_6500: mean spread 1, mean mid 0.5, trades/day near 91-98.
\end{Verbatim}

\edaplot{voucher_mid_summary_grid_normalized.png}{Normalized voucher mid-price complex.}
\edaplot{voucher_mid_VEV_5200.png}{Representative near-ATM voucher mid: VEV 5200.}
\edaplot{voucher_mid_VEV_5300.png}{Representative near-ATM voucher mid: VEV 5300.}
\edaplot{voucher_mid_VEV_5500.png}{Representative OTM active voucher mid: VEV 5500.}
\edaplot{voucher_mid_VEV_6000.png}{Pinned floor-price voucher mid: VEV 6000.}
\edaplot{voucher_spread_VEV_4000.png}{Deep ITM voucher spread distribution.}
\edaplot{voucher_spread_VEV_5300.png}{Near-ATM voucher spread distribution.}
\edaplot{voucher_spread_VEV_6500.png}{Deep OTM floor-price voucher spread distribution.}

\section{Options Valuation, IV, Greeks, And Arbitrage}

\subsection*{Black-Scholes And IV Findings}

\begin{itemize}
  \item The longer EDA assumes Black-Scholes calls with zero risk-free rate and \code{T_years = TTE_days / 365}.
  \item Historical TTE mapping is day 0 = 8 days, day 1 = 7 days, day 2 = 6 days; live Round 3 is expected to have 5 days to expiry.
  \item ATM IV is selected as the closest valid-IV voucher at each timestamp; this is almost always \code{VEV_5200} or \code{VEV_5300}.
  \item ATM IV has mean 0.232 on day 0, 0.237 on day 1, and 0.235 on day 2.
  \item RV500 has mean 0.410, 0.412, and 0.414 on days 0, 1, and 2.
  \item Therefore IV is persistently below realized volatility by about 17.5 to 18.0 vol points.
  \item The IV smile is smooth: quadratic smile fit mean \(R^2\) is about 0.993, 0.989, and 0.987 across the three days.
  \item \code{VEV_6000} and \code{VEV_6500} have high implied vols but they are mostly artifacts of the 0/1 quote floor and should be treated carefully.
\end{itemize}

\begin{Verbatim}[fontsize=\scriptsize]
ATM IV versus RV500:
day  mean ATM IV  mean RV500  mean IV - RV
0    0.232175     0.409696   -0.177877
1    0.237291     0.412265   -0.175404
2    0.234733     0.413823   -0.179891

IV status counts:
ok                   251785
zero_time_value       36414
below_lower_bound     11801

Main invalid-IV sources:
VEV_4000: below_lower_bound 2986, zero_time_value 24028.
VEV_4500: below_lower_bound 8812, zero_time_value 12365.
VEV_5000: below_lower_bound 3, zero_time_value 21.
\end{Verbatim}

\edaplot{critical_atm_iv_vs_velvetfruit_rv.png}{Critical chart: ATM IV versus underlying rolling realized volatility.}
\edaplot{voucher_atm_iv_across_days.png}{ATM IV proxy across historical days.}
\edaplot{voucher_atm_iv_ema_bands.png}{ATM IV with EMA bands.}
\edaplot{voucher_iv_timeseries_all.png}{Per-voucher implied volatility over time.}
\edaplot{voucher_iv_smile_day0.png}{IV smile snapshots, day 0.}
\edaplot{voucher_iv_smile_day1.png}{IV smile snapshots, day 1.}
\edaplot{voucher_iv_smile_day2.png}{IV smile snapshots, day 2.}
\edaplot{voucher_iv_smile_smoothness.png}{Quadratic IV-smile smoothness over time.}
\edaplot{voucher_iv_term_structure.png}{Cross-day IV term-structure proxy.}

\subsection*{Greeks, Intrinsic Value, And Gamma-Scalp Proxy}

\begin{itemize}
  \item Median-sample delta declines monotonically from near 1.0 in deep ITM strikes to near zero in deep OTM strikes.
  \item Gamma is concentrated around \code{VEV_5200} and \code{VEV_5300}, with \code{VEV_5100} and \code{VEV_5400} also relevant.
  \item \code{VEV_4000} and \code{VEV_4500} have almost no usable gamma despite high trade count or high notional price.
  \item Intrinsic/extrinsic decomposition confirms that \code{VEV_4000} and \code{VEV_4500} are almost pure intrinsic value.
  \item Extrinsic value decays with TTE in \code{VEV_5000}, \code{VEV_5100}, \code{VEV_5200}, and OTM strikes.
  \item Gamma-scalp proxy is positive for \code{VEV_4500} through \code{VEV_6000}, with the largest proxy in \code{VEV_5200} and \code{VEV_5300}.
  \item The proxy is negative for \code{VEV_4000} and \code{VEV_6500}, where theta-like cost dominates realized gamma benefit.
\end{itemize}

\begin{Verbatim}[fontsize=\scriptsize]
Median greeks:
VEV_4000 delta 0.9970, gamma 0.000017
VEV_5000 delta 0.9286, gamma 0.000796
VEV_5100 delta 0.8174, gamma 0.001619
VEV_5200 delta 0.6076, gamma 0.002261
VEV_5300 delta 0.3794, gamma 0.002192
VEV_5400 delta 0.1745, gamma 0.001588
VEV_5500 delta 0.0732, gamma 0.000815

Gamma-scalp proxy, 3-day net per 1 option unit:
VEV_5200  +29.18
VEV_5300  +28.92
VEV_5400  +22.48
VEV_5100  +20.29
VEV_5500  +11.07
VEV_5000   +9.50
VEV_6500   -0.52
VEV_4000   -0.12
\end{Verbatim}

\edaplot{voucher_delta_vs_strike.png}{Median-sample Black-Scholes delta by strike.}
\edaplot{voucher_gamma_vs_strike.png}{Median-sample Black-Scholes gamma by strike.}
\edaplot{voucher_delta_gamma_timeseries.png}{Delta and gamma time series by voucher.}
\edaplot{voucher_intrinsic_extrinsic.png}{Intrinsic and extrinsic decomposition for selected vouchers.}
\edaplot{gamma_scalp_pnl_summary.png}{Gamma-scalp PnL proxy summary.}

\subsection*{No-Arbitrage Diagnostics}

\begin{itemize}
  \item Lower-bound violations dominate the no-arbitrage table: 11,801 total flags.
  \item There are no upper-bound violations.
  \item Lower-bound flags are concentrated in \code{VEV_4500} and \code{VEV_4000}.
  \item There are 9 convexity flags involving \code{VEV_5400}, \code{VEV_5500}, and \code{VEV_6000}.
  \item Average lower-bound flag lifetimes are short, around 1.1 to 1.5 timestamps.
  \item Practical strategy must require executable bid/ask edge, persistence, and inventory feasibility before treating these as trades.
\end{itemize}

\begin{Verbatim}[fontsize=\scriptsize]
Positive no-arbitrage flags:
day  voucher/group                   type          count  avg_magnitude  avg_lifetime
0    VEV_4000                        lower_bound     939      1.2268        1.097
0    VEV_4500                        lower_bound    3045      0.6842        1.485
0    VEV_5400,VEV_5500,VEV_6000      convexity         9      0.5000        1.000
1    VEV_4000                        lower_bound     984      1.2175        1.147
1    VEV_4500                        lower_bound    3061      0.6900        1.484
2    VEV_4000                        lower_bound    1063      1.1834        1.167
2    VEV_4500                        lower_bound    2706      0.7308        1.420
2    VEV_5000                        lower_bound       3      0.5000        1.000
\end{Verbatim}

\section{Lead-Lag, Residuals, And Trade-Against-Fair-Value}

\begin{itemize}
  \item Lead-lag correlations peak at lag 0 for all vouchers with meaningful price variation.
  \item Peak correlations are highest for \code{VEV_5100}, \code{VEV_5000}, and \code{VEV_5200}.
  \item \code{VEV_6000} and \code{VEV_6500} have undefined lead-lag correlations because their mid prices are pinned.
  \item Delta-hedged residual standard deviation increases as the option becomes more nonlinear and illiquid.
  \item Delta-hedged residual autocorrelation is negative across all active vouchers, strongest around \code{VEV_5400} and \code{VEV_5500}; this suggests short-horizon residual mean reversion rather than momentum.
  \item Voucher trades versus Black-Scholes fair value show that \code{VEV_5200} and higher active strikes usually trade below model fair value, while the few valid \code{VEV_4000}, \code{VEV_5000}, and \code{VEV_5100} trades are above fair value.
  \item Maker/taker fill breakdown shows aggressive selling pressure in \code{VEV_5200} through \code{VEV_6500}: almost all prints occur at the bid.
\end{itemize}

\begin{Verbatim}[fontsize=\scriptsize]
Lead-lag summary:
voucher   best_lag  peak_corr  residual_std  residual_acf1
VEV_5000     0       0.7524       0.00371       -0.1045
VEV_5100     0       0.7627       0.00502       -0.1025
VEV_5200     0       0.7193       0.00713       -0.1382
VEV_5300     0       0.6123       0.01081       -0.2217
VEV_5400     0       0.5252       0.01770       -0.2508
VEV_5500     0       0.3338       0.02855       -0.2425

Trade price minus BS FV:
VEV_4000   +1.446, but only 46 valid-IV trades.
VEV_5200   -0.917
VEV_5300   -0.963
VEV_5400   -0.629
VEV_5500   -0.547
VEV_6000   -0.500
VEV_6500   -0.500
\end{Verbatim}

\edaplot{leadlag_corr_VEV_5100.png}{Lead-lag correlation for VEV 5100.}
\edaplot{leadlag_corr_VEV_5200.png}{Lead-lag correlation for VEV 5200.}
\edaplot{leadlag_corr_VEV_5300.png}{Lead-lag correlation for VEV 5300.}
\edaplot{delta_hedged_residual_hist_VEV_5200.png}{Delta-hedged residual distribution for VEV 5200.}
\edaplot{delta_hedged_residual_hist_VEV_5300.png}{Delta-hedged residual distribution for VEV 5300.}
\edaplot{voucher_trade_vs_bs_fv_hist.png}{Voucher trade price minus Black-Scholes fair value.}

\section{Consistency Check Between The Two EDAs}

\begin{itemize}
  \item Target product counts are consistent. The notebook's day-by-day trade counts match the longer EDA's \code{data_integrity.csv} and \code{voucher_trade_sizes.csv}.
  \item Data sources are effectively equivalent for target products: the notebook reads combined files, while the long EDA reads per-day files.
  \item The main denominator inconsistency is the notebook's product-frequency section: it prints frequencies after filtering out \code{HYDROGEL_PACK}, while the full dataset has 12 products and 360,000 rows.
  \item Missing depth is handled differently. The notebook converts missing volumes to zero before computing depth and anomaly states; the long EDA preserves missing price levels and explicitly measures deeper-book availability.
  \item The notebook is a structural/microstructure EDA, while the long EDA is an option-complex EDA. There is no conflict on IV, greeks, no-arbitrage, or gamma findings because the notebook does not compute those quantities.
  \item The notebook's depletion result for \code{VEV_6500} is consistent with the long EDA's finding that deep OTM vouchers trade at bid 0 with a bid/ask of 0/1.
  \item The notebook's advanced metric \code{deep_vamp} is useful descriptively, but it should not be compared directly to the longer EDA's Black-Scholes fair values or delta-hedged residuals.
\end{itemize}

\section{Code Audit Notes}

\begin{itemize}
  \item No fatal syntax error was found in \code{ve_vev_eda.py}; the file parses successfully.
  \item A full runtime check using \code{--output-dir /tmp/vev_eda_check} completed successfully without overwriting the existing output directory.
  \item The runtime produced a KPSS interpolation warning. This is not fatal; it means the KPSS statistic is outside the tabulated p-value range.
  \item The requested path \code{phase2/round3/algo/analysis/ve_vev_eda.py} does not exist in the workspace; the actual script is \code{phase2/round3/algo/ve_vev_eda.py}.
  \item \code{ensure_output()} deletes all existing files in the selected charts and tables directories. This is intended for regeneration but dangerous if pointed at a non-temporary directory.
  \item The scalar \code{bs_greeks()} divides vega by 100, while \code{bs_greeks_vec()} does not. As a result, \code{voucher_greeks.csv} and \code{voucher_greeks_timeseries.csv} use different vega units.
  \item ACF, variance-ratio, Hurst, and stationarity diagnostics are partly computed over the pooled three-day series, so day-boundary jumps can influence the tests.
  \item The regime label "trending" conflicts with variance ratios below 1.0 and negative return autocorrelation. This should be interpreted as "price level persistent with local mean-reverting ticks", not as simple tradable momentum.
  \item The IV term structure plot is a historical cross-day proxy, not a true same-time term structure, because TTE changes with day and the underlying path changes too.
  \item Quote-change rates count the first row of each day as a change because \code{diff() != 0} treats the initial NaN as changed; the bias is tiny but present.
  \item \code{velvetfruit_bot_volume_stats.csv} writes a MultiIndex aggregation without flattening headers, producing an awkward first row containing \code{mean/std/median}.
  \item Section numbering in printed runtime messages says \code{1/6} through \code{5/6}, then \code{6/7} and \code{7/7}; this is cosmetic.
\end{itemize}

\section{Initial Trading Implications}

\begin{itemize}
  \item Treat \code{VELVETFRUIT_EXTRACT} as the hedge instrument. Its spread and depth are stable enough for controlled delta hedging, but crossing costs about 2.5 ticks versus mid.
  \item The largest systematic signal is IV below RV. The natural hypothesis is to own gamma near ATM and delta-hedge, especially in \code{VEV_5200}, \code{VEV_5300}, \code{VEV_5400}, and \code{VEV_5100}.
  \item Because voucher spreads and fills are sparse, the strategy should be passive-first in vouchers and selective-taker in the underlying when hedge error becomes too large.
  \item Deep ITM vouchers are mostly intrinsic value and lower-bound flags are not automatically tradable after spreads.
  \item Deep OTM floor vouchers should be handled as special discrete-state instruments, not continuous Black-Scholes options.
  \item A residual mean-reversion overlay can be tested: quote options when delta-hedged residuals are extreme relative to local bands, then hedge with Velvetfruit.
  \item Any production strategy should combine surface fair value, IV/RV regime, expected fill probability, inventory Greeks, and spread-adjusted expected PnL.
\end{itemize}

\section{100 Research Questions For Further Analysis}

\begin{enumerate}[label=Q\arabic*.,leftmargin=*]
  \item How stable is the IV-minus-RV gap if RV is computed with windows other than 500 ticks?
  \item Does the IV-minus-RV gap predict future delta-hedged option PnL after spread and hedge costs?
  \item Which voucher has the highest net expected gamma-scalp PnL after realistic fills, not midpoint marks?
  \item Does the gap between ATM IV and RV mean-revert intraday?
  \item Does the ATM IV EMA band signal produce profitable entry and exit points?
  \item Are \code{VEV_5200} and \code{VEV_5300} always the best ATM proxies, or does the closest strike switch in predictive ways?
  \item How much of the gamma-scalp proxy survives if the underlying hedge is executed at bid/ask instead of mid?
  \item What hedge threshold minimizes total error plus transaction costs for a long-gamma option book?
  \item How often can we obtain voucher fills passively at favorable model prices?
  \item Which side of each voucher book receives more aggressive flow conditional on time of day?
  \item Does underlying buy-flow toxicity predict voucher mid moves after controlling for delta?
  \item Does underlying order-book imbalance become predictive only in high-volume or high-volatility regimes?
  \item Is the underlying lag-1 return autocorrelation tradable after spread costs?
  \item Are underlying price moves larger after clustered trade-arrival buckets?
  \item Does the intraday buy fraction predict the next 1000-timestamp underlying return?
  \item Are underlying spreads wider before volatility bursts, or merely contemporaneous with RV?
  \item Does total depth predict near-future volatility or trade direction?
  \item Do quote offsets from mid change before large underlying moves?
  \item Is the bid/ask quote-change rate informative about the next mid-price move?
  \item Do outside-top-book underlying trades have special predictive content?
  \item Can \code{VELVETFRUIT_EXTRACT} be profitably market-made using only mid reversion and stable spread?
  \item How does optimal underlying hedge size vary with option gamma and current underlying spread?
  \item Does buying at the underlying bid after sell pressure outperform passive random quoting?
  \item Are day 1 and day 2 positive drifts persistent enough to justify directional delta?
  \item Does the underlying show regime-switching between drift, range, and shock states?
  \item Which voucher's delta-hedged residual has the strongest post-cost mean reversion?
  \item Are residual bands stable by day, or must they be re-estimated intraday?
  \item Do residual shocks in one voucher predict residual shocks in nearby strikes?
  \item Does the voucher complex move as a low-dimensional IV surface factor model?
  \item Can PCA factors of IV surface predict future option returns better than per-strike IV levels?
  \item Are no-arbitrage lower-bound flags executable at bid/ask, or only midpoint artifacts?
  \item How long must a no-arbitrage flag persist before it is worth crossing a spread?
  \item Are convexity violations around \code{VEV_5400}/\code{VEV_5500}/\code{VEV_6000} repeatable outside day 0?
  \item Can monotonic and convex smoothing of option mids improve fair-value quotes?
  \item Does a constrained IV smile outperform raw mid IV for predicting future trades?
  \item Do deep ITM vouchers behave more like stock-plus-cash than options?
  \item Is \code{VEV_4000} useful as a synthetic underlying hedge when the underlying spread is unfavorable?
  \item Can a box or vertical-spread portfolio lock in edge after accounting for all bid/ask spreads?
  \item Does the option complex imply a forward price that differs from Velvetfruit spot?
  \item Are integer-price lower-bound violations caused by rounding rules rather than mispricing?
  \item Is there profitable edge in posting bid 0 for \code{VEV_6000} and \code{VEV_6500}?
  \item Is there profitable edge in posting ask 1 for floor-price OTM vouchers, or is tail risk too high?
  \item How should tail risk be modeled for short \code{VEV_6000}/\code{VEV_6500} positions?
  \item Do floor-price vouchers ever move away from 0/1 under different live paths?
  \item Is the empirical IV of floor-price vouchers meaningful, or should they be modeled as binary lottery tickets?
  \item How does TTE decline from 6 historical days to 5 live days change the ATM fair value?
  \item Can we extrapolate IV surfaces safely from TTE 8/7/6 to live TTE 5?
  \item Does theta decay estimated from historical day transitions match Black-Scholes theta?
  \item Is the observed drop in extrinsic value across days mostly theta or underlying-path effect?
  \item Which strikes have the most stable extrinsic value after conditioning on moneyness?
  \item Are voucher trade sizes informative about informed flow?
  \item Does a large voucher trade at bid predict lower future IV?
  \item Does a voucher trade at ask predict higher future IV or just temporary demand?
  \item How much model edge is needed before crossing a voucher spread is profitable?
  \item What passive quote distance maximizes expected fill-adjusted PnL for each strike?
  \item Can fill probability be estimated from quote position, spread, and recent trade direction?
  \item Are voucher books dominated by stable bots with predictable posted volumes?
  \item Does bid volume or ask volume asymmetry in vouchers predict the next voucher trade side?
  \item Does voucher quote-change rate decline monotonically with OTM distance because prices are pinned?
  \item Can stable volume patterns identify which bot is quoting each strike?
  \item Are the buyer/seller columns truly empty in all live-style historical data?
  \item How sensitive are Black-Scholes fair values to using calendar-day versus trading-day annualization?
  \item Should volatility be annualized using 365 days, 7-day option expiry, or Prosperity-specific tick time?
  \item Does realized volatility measured on log returns overstate volatility because of bid/ask bounce?
  \item How does RV change after filtering zero changes or using mid-price smoothing?
  \item Does using microprice instead of mid improve underlying volatility and delta estimates?
  \item Is deep VAMP more predictive of future underlying moves than top-level imbalance?
  \item Does DOM-mid deviation predict future reversion to top-of-book mid?
  \item Can underlying fair value be improved with Kalman filtering of noisy mid prices?
  \item How different are option deltas if fair spot is filtered rather than raw mid?
  \item Does delta hedging with filtered spot reduce residual variance?
  \item Does a smile-fitted IV produce lower delta-hedged residual variance than raw per-strike IV?
  \item Which residual definition best predicts actual trade PnL: price residual, return residual, or dollar hedge error?
  \item Can residual autocorrelation be exploited with passive quotes only?
  \item Are residuals heteroskedastic by strike and time to expiry?
  \item Does residual volatility rise near day boundaries or late-session windows?
  \item Is there a consistent intraday seasonality in ATM IV?
  \item Does the IV smile steepen or flatten as underlying spot moves?
  \item Is skew/smile curvature predictive of future spot direction?
  \item Can a straddle or strangle portfolio isolate the IV-vs-RV edge better than single options?
  \item Which option portfolio minimizes delta while maximizing gamma under position limits?
  \item What is the optimal allocation of the 300-lot voucher limits across strikes?
  \item What is the optimal \code{VELVETFRUIT_EXTRACT} hedge inventory under a portfolio of voucher positions?
  \item How should inventory risk be penalized near the end of the simulation?
  \item Does forced mark-to-market at final timestamp favor closing or holding option inventory?
  \item What worst-case underlying move breaks a short-OTM voucher strategy?
  \item What worst-case spread widening breaks a long-gamma strategy?
  \item Are model signals robust if one historical day is held out?
  \item Which features survive walk-forward validation from day 0 to day 1 to day 2?
  \item Does a live TTE=5 strategy trained on TTE=6 require parameter shrinkage?
  \item How should we combine directional underlying alpha with option gamma alpha?
  \item Can we synthesize a better underlying exposure using deep ITM vouchers plus cash?
  \item What are the PnL effects of integer order prices and tick-size rounding?
  \item How often does rounding flip a theoretical edge from positive to negative?
  \item What is the best quoting grid for each voucher after fair-value rounding?
  \item Which signals remain profitable after adverse-selection penalties from fills?
  \item Can a joint simulator replay historical order books and test passive fills realistically?
  \item Which EDA signal is most likely to overfit only three historical days?
  \item What minimal strategy captures most of the IV/RV edge with the least operational risk?
  \item What live-regime guardrail should respond to RV collapse, IV spikes, or hedge PnL deviation?
\end{enumerate}

\section{30 Candidate Models And Validation Statistics}

\begin{enumerate}[label=M\arabic*.,leftmargin=*]
  \item \textbf{Random walk with drift for underlying mid.} Validate with out-of-sample log-likelihood, drift t-statistics, ADF/KPSS on residual increments, and variance-ratio tests.
  \item \textbf{ARMA/ARIMA on underlying returns.} Validate with AIC/BIC, Ljung-Box residual tests, residual ACF, walk-forward RMSE, and directional hit rate.
  \item \textbf{ARFIMA or fractional-noise model for price persistence.} Validate with Hurst consistency, Whittle likelihood, residual long-memory tests, and out-of-sample forecast loss.
  \item \textbf{Ornstein-Uhlenbeck model for microstructure return reversion.} Validate with half-life stability, residual normality, likelihood ratio versus random walk, and post-cost mean-reversion PnL.
  \item \textbf{Regime-switching AR model for drift/range/shock states.} Validate with hidden-state persistence, out-of-sample log score, transition stability by day, and strategy PnL by inferred regime.
  \item \textbf{GARCH(1,1) for underlying one-tick returns.} Validate with QLIKE volatility loss, ARCH-LM residual tests, standardized residual kurtosis, and volatility forecast MSE.
  \item \textbf{EGARCH or GJR-GARCH for asymmetric volatility.} Validate with leverage-term significance, AIC/BIC improvement over GARCH, VaR exception counts, and residual diagnostics.
  \item \textbf{HAR-RV model using RV100/RV500/RV2000.} Validate with out-of-sample RV forecast RMSE/MAE, QLIKE, Diebold-Mariano tests versus EWMA, and coefficient stability.
  \item \textbf{EWMA/RiskMetrics volatility model.} Validate lambda via walk-forward likelihood, RV forecast error, and sensitivity of option fair values to lambda.
  \item \textbf{Kalman local-level model for filtered underlying fair value.} Validate with one-step prediction error, filtered-versus-raw hedge residual variance, and likelihood.
  \item \textbf{Order-book imbalance linear model for forward returns.} Validate with t-statistics, adjusted \(R^2\), robust standard errors, sign hit rate, and post-spread PnL.
  \item \textbf{Nonlinear tree or spline model for imbalance and depth features.} Validate with cross-validated \(R^2\), permutation importance, calibration by feature bucket, and walk-forward PnL.
  \item \textbf{Ordered logit/probit for next tick move direction.} Validate with log loss, Brier score, confusion matrix, calibration curves, and economic value after costs.
  \item \textbf{Poisson or negative-binomial model for trade arrivals.} Validate with dispersion index, likelihood ratio test, Pearson residuals, and intraday bucket calibration.
  \item \textbf{Hawkes process for clustered trades.} Validate with time-rescaling KS test, branching ratio stability, likelihood improvement over Poisson, and burst-period PnL attribution.
  \item \textbf{Autoregressive conditional duration model for interarrival times.} Validate with standardized duration diagnostics, AIC/BIC, and forecasted fill-risk calibration.
  \item \textbf{Black-Scholes fair-value model with empirical IV.} Validate with trade-minus-fair residuals, delta-hedged residual variance, IV status coverage, and post-cost option PnL.
  \item \textbf{Constrained quadratic IV smile by timestamp.} Validate with smile \(R^2\), no-arbitrage violation count, out-of-sample trade-price error, and residual smoothness.
  \item \textbf{SVI-style parametric implied-volatility surface.} Validate with weighted IV RMSE, convexity/no-arbitrage checks, parameter stability, and bid/ask fit rate.
  \item \textbf{SABR-inspired smile approximation.} Validate with strike-weighted IV error, curvature fit around ATM, parameter plausibility, and delta/gamma hedge error.
  \item \textbf{PCA/factor model of IV surface changes.} Validate with explained variance, factor autocorrelation, out-of-sample reconstruction error, and factor-PnL relationship.
  \item \textbf{Bayesian hierarchical IV model by strike and day.} Validate posterior predictive checks, credible-interval coverage, leave-one-day-out likelihood, and shrinkage robustness.
  \item \textbf{Delta-hedged residual AR(1) model by voucher.} Validate residual half-life, residual ACF, z-score calibration, mean-reversion PnL, and turnover-adjusted Sharpe.
  \item \textbf{Gamma-scalping PnL attribution model.} Validate realized \(0.5\Gamma dS^2 + \Theta dt\) attribution, hedge-cost drag, net PnL by strike, and sensitivity to hedge frequency.
  \item \textbf{No-arbitrage persistence survival model.} Validate hazard calibration, precision/recall for persistent flags, spread-adjusted opportunity PnL, and false-positive rate.
  \item \textbf{Constrained cross-sectional option-pricing regression.} Validate monotonicity, convexity, bid/ask inclusion rate, residual RMSE, and executable arbitrage count.
  \item \textbf{Voucher fill-probability logistic model.} Validate ROC-AUC, precision at top quote candidates, calibration curves, and expected fill-adjusted PnL.
  \item \textbf{Avellaneda-Stoikov style market-making model with option Greeks.} Validate inventory distribution, spread capture, adverse-selection loss, drawdown, and risk-adjusted PnL.
  \item \textbf{Linear-quadratic inventory control for delta/gamma constrained portfolios.} Validate terminal inventory penalty, hedge turnover, realized PnL distribution, and stress-test loss.
  \item \textbf{Walk-forward ensemble combining IV/RV, residual, flow, and inventory signals.} Validate day-held-out PnL, feature ablation, turnover-adjusted Sharpe, maximum drawdown, and live TTE=5 sensitivity.
\end{enumerate}

\section{Appendix: Additional Plot Panels}

The following panels include the main plot families generated by the longer EDA. They are deliberately kept as visual appendices so the core research narrative stays readable.

\begin{figure}[H]
\centering
\smallplot{voucher_mid_VEV_4000.png}
\smallplot{voucher_mid_VEV_4500.png}
\smallplot{voucher_mid_VEV_5000.png}
\smallplot{voucher_mid_VEV_5100.png}
\smallplot{voucher_mid_VEV_5200.png}
\smallplot{voucher_mid_VEV_5300.png}
\smallplot{voucher_mid_VEV_5400.png}
\smallplot{voucher_mid_VEV_5500.png}
\smallplot{voucher_mid_VEV_6000.png}
\smallplot{voucher_mid_VEV_6500.png}
\caption{Per-voucher mid-price time series.}
\end{figure}

\begin{figure}[H]
\centering
\smallplot{voucher_spread_VEV_4000.png}
\smallplot{voucher_spread_VEV_4500.png}
\smallplot{voucher_spread_VEV_5000.png}
\smallplot{voucher_spread_VEV_5100.png}
\smallplot{voucher_spread_VEV_5200.png}
\smallplot{voucher_spread_VEV_5300.png}
\smallplot{voucher_spread_VEV_5400.png}
\smallplot{voucher_spread_VEV_5500.png}
\smallplot{voucher_spread_VEV_6000.png}
\smallplot{voucher_spread_VEV_6500.png}
\caption{Per-voucher spread distributions.}
\end{figure}

\begin{figure}[H]
\centering
\smallplot{leadlag_corr_VEV_4000.png}
\smallplot{leadlag_corr_VEV_4500.png}
\smallplot{leadlag_corr_VEV_5000.png}
\smallplot{leadlag_corr_VEV_5100.png}
\smallplot{leadlag_corr_VEV_5200.png}
\smallplot{leadlag_corr_VEV_5300.png}
\smallplot{leadlag_corr_VEV_5400.png}
\smallplot{leadlag_corr_VEV_5500.png}
\smallplot{leadlag_corr_VEV_6000.png}
\smallplot{leadlag_corr_VEV_6500.png}
\caption{Lead-lag correlation panels by voucher.}
\end{figure}

\begin{figure}[H]
\centering
\smallplot{gamma_scalp_pnl_VEV_4000.png}
\smallplot{gamma_scalp_pnl_VEV_4500.png}
\smallplot{gamma_scalp_pnl_VEV_5000.png}
\smallplot{gamma_scalp_pnl_VEV_5100.png}
\smallplot{gamma_scalp_pnl_VEV_5200.png}
\smallplot{gamma_scalp_pnl_VEV_5300.png}
\smallplot{gamma_scalp_pnl_VEV_5400.png}
\smallplot{gamma_scalp_pnl_VEV_5500.png}
\smallplot{gamma_scalp_pnl_VEV_6000.png}
\smallplot{gamma_scalp_pnl_VEV_6500.png}
\caption{Gamma-scalp proxy time series by voucher.}
\end{figure}

\begin{figure}[H]
\centering
\smallplot{delta_hedged_residual_series_VEV_5000.png}
\smallplot{delta_hedged_residual_hist_VEV_5000.png}
\smallplot{delta_hedged_residual_series_VEV_5100.png}
\smallplot{delta_hedged_residual_hist_VEV_5100.png}
\smallplot{delta_hedged_residual_series_VEV_5200.png}
\smallplot{delta_hedged_residual_hist_VEV_5200.png}
\smallplot{delta_hedged_residual_series_VEV_5300.png}
\smallplot{delta_hedged_residual_hist_VEV_5300.png}
\smallplot{delta_hedged_residual_series_VEV_5400.png}
\smallplot{delta_hedged_residual_hist_VEV_5400.png}
\caption{Representative delta-hedged residual series and histograms.}
\end{figure}

\end{document}
